\documentclass[11pt]{article}
\usepackage{preamble}
\renewcommand{\answer}[1]{}

\begin{document}

\ladderheading{AP Statistics \textperiodcentered\ Unit 5 \textperiodcentered\ Topic 5.5 \textperiodcentered\ B: 2027-02-12 \textperiodcentered\ E: 2027-04-09}{Which Regression Line Fits?}

\focusbanner{TODAY'S PROBLEM}

Suppose a solar-farm analyst uses hours of direct sunshine, $x$, to predict daily generation, $y$, in megawatt-hours. Summary statistics for 12 observed days are shown.\par\smallskip \textbf{Quinn:} ``Using $b=r(s_y/s_x)$ gives $\hat y=5.6+3.2x$. It passes through $(7,28)$, which by itself proves it is the LSRL. Since $r=0.80$, the model explains 80 percent of the variation.''\par \textbf{Rowan:} ``Using $b=r(s_x/s_y)$ gives $\hat y=26.6+0.2x$. It also passes through $(7,28)$, so that line is equally eligible. Explained variation is 64 percent, and a zero-hour intercept needs caution because 0 is outside the observed range.''\par
\begin{center}
\textbf{Solar-farm summary (12 observed days)}\par\smallskip
\resizebox{0.98\linewidth}{!}{%
\begin{tabular}{|c|c|c|c|c|c|}
\hline
$x$ range & $\bar x$ & $\bar y$ & $s_x$ & $s_y$ & $r$ \\ \hline
4--10 h & 7 h & 28 MWh & 1.5 h & 6 MWh & 0.80 \\ \hline
\end{tabular}}
\end{center}
\begin{firsttakebox}{FIRST TAKE --- before the video (5 min)}
Which clauses in Quinn's and Rowan's analyses look defensible? Cite one formula or line property.
\workspace{0.65in}
\end{firsttakebox}

\begin{callout}{\IconBook\ MAKE THE REGRESSION PROPERTIES AGREE (keep on the page)}{calloutgreen}
The least-squares regression line (LSRL) minimizes the sum of squared residuals and contains
$(\bar x,\bar y)$. Its slope is $b=r(s_y/s_x)$, its intercept is
$a=\bar y-b\bar x$, and $r^2$ is the proportion of response variation explained by the
linear model. Interpret slope as a predicted response change per unit of $x$. The intercept
predicts at $x=0$, which may be outside the observed range and unsupported in context.
\end{callout}

\newpage
\textbf{Finish after the video:}\par\smallskip

\dokbadge{1}~\textbf{(a)} State what the LSRL minimizes, the point it contains, and the contextual meanings of slope and intercept.\par
\workspace{0.55in}

\dokbadge{2}~\textbf{(b)} Calculate the LSRL slope and intercept, $r^2$, and the prediction at 8 sunshine hours. Verify the mean-point property.\par
\workspace{1.0in}

\dokbadge{3}~\textbf{(c)} Evaluate both analyses and write one warranted synthesis. Coordinate the slope ratio and units, least-squares and mean-point properties, coefficient meanings, and the observed domain. Compare both lines at 8 hours and quantify the consequences of using the rejected slope and using $r$ instead of $r^2$.\par
\workspace{1.45in}

\begin{summaryexitbox}{EXIT --- turn this in}
One thing the video changed about my first take: \hrulefill\par
\workspace{0.4in}
\end{summaryexitbox}

\end{document}
